\section{Related Work}
\label{sec:related}

\para{Temporal degradation in language models.}
Several recent studies establish that LLM performance degrades with temporal distance from training data. \citet{li2024outdated} introduce FreshBench and show that all tested models exhibit positive temporal bias on news text, with finance and business categories degrading fastest. They also identify a ``nostalgia bias''---models peak at knowledge from 3--5 years before their cutoff, not at the cutoff itself. \citet{nylund2023time} demonstrate that temporal misalignment degrades performance linearly at roughly 7--10 F1 points per year, and show that temporal information is encoded primarily in feed-forward layers rather than embeddings. \citet{pelcra2025llmlagbench} find that declared knowledge cutoffs frequently diverge from empirical ones by 1--2 years. Our work extends these findings by measuring not just \emph{that} degradation occurs, but \emph{how} the fine-tuning learning process itself differs across domains---revealing that convergence dynamics, not just final performance, vary systematically.

\para{Knowledge injection via fine-tuning.}
\citet{ovadia2024finetuning} directly compare fine-tuning and retrieval-augmented generation (RAG) for injecting post-cutoff knowledge, finding that RAG consistently outperforms fine-tuning and that fine-tuning on new data can even decrease accuracy. \citet{longpre2023pretrainer} show that temporal distribution mismatch between pretraining and evaluation data causes degradation that fine-tuning cannot overcome. \citet{zhao2024set} introduce temporal alignment---shifting a model's internal ``knowledge clock'' via fine-tuning on time-indexed QA data---and find that while alignment activates existing suppressed knowledge, 61.5\% of recently-changed facts remain wrong even after alignment. Building on these results, we focus on \emph{convergence dynamics} rather than end-state accuracy, showing that the learning process itself reveals whether a model is absorbing facts or merely learning surface patterns.

\para{Fundamental limitations of autoregressive learning.}
\citet{berglund2024reversal} demonstrate the reversal curse: models trained on ``A is B'' completely fail to learn ``B is A,'' even at 175B parameters. This asymmetry persists across all model sizes and data augmentation strategies, suggesting a structural limitation of autoregressive training for knowledge acquisition. \citet{shumailov2023curse} show that training on model-generated data leads to progressive loss of distributional tails (model collapse), implying that older models may be structurally unable to simulate the full diversity of future text. These architectural constraints compound the temporal knowledge problem: even when fine-tuning exposes a model to new facts, the learning signal may be too weak or too directional to produce robust knowledge representations.

\para{Continual learning and temporal adaptation.}
The broader continual learning literature addresses how models can acquire new knowledge without forgetting old knowledge \citep{jang2022temporalwiki}. Our work differs in focus: rather than developing new methods for temporal adaptation, we characterize \emph{which types of knowledge} resist adaptation and why. This diagnostic perspective complements method-oriented work by identifying where fine-tuning is a viable strategy and where alternative approaches like RAG \citep{lewis2020rag} are necessary.
